% The report's packages, styles and shared macros. Every section uses these; add what several
% sections need here, and what one section needs at that section's top (with \providecommand).

\usepackage[utf8]{inputenc}
\usepackage[T1]{fontenc}
\usepackage{lmodern}
\usepackage{microtype}
\usepackage{amsmath,amssymb,bm,mathtools}
\usepackage[a4paper,margin=2.5cm]{geometry}
\usepackage{graphicx}
\usepackage{booktabs,tabularx,longtable,multirow}
\usepackage{enumitem}
\usepackage{siunitx}
\usepackage{xcolor}
% The report's colours: three that stay apart for every kind of colour vision, in a fixed
% order (the solver's curve is always blue), five more for figures with many curves; the
% reference (an exact or a published answer) is drawn in black.

\definecolor{vrblue}{HTML}{257DE0}
\definecolor{vrorange}{HTML}{F47230}
\definecolor{vraqua}{HTML}{00B894}
\definecolor{vryellow}{HTML}{F2A600}
\definecolor{vrmagenta}{HTML}{EE7FAC}
\definecolor{vrgreen}{HTML}{008A17}
\definecolor{vrviolet}{HTML}{6A4CE2}
\definecolor{vrred}{HTML}{EC4647}
\definecolor{vrlink}{HTML}{0F4A96}
\definecolor{vrgrey}{HTML}{666666}


\usepackage{tikz}
\usetikzlibrary{arrows.meta,calc,positioning,shapes.geometric,patterns,decorations.pathreplacing}
\usepackage{pgfplots,pgfplotstable}
\usepgfplotslibrary{groupplots,fillbetween}
\usepackage{listings}
\usepackage[most]{tcolorbox}
\usepackage[font=small,labelfont=bf]{caption}
\usepackage{subcaption}
\usepackage[section]{placeins}
\usepackage[round,sort&compress]{natbib}
\PassOptionsToPackage{obeyspaces,spaces}{url}   % \code keeps its spaces
\usepackage[colorlinks=true,linkcolor=vrlink,citecolor=vrlink,urlcolor=vrlink]{hyperref}
\usepackage[capitalise,noabbrev]{cleveref}
% The bibliography style writes \doi{...}: a link.
\newcommand{\doi}[1]{\href{https://doi.org/#1}{\nolinkurl{doi:#1}}}

% ---------------------------------------------------------------------------------------------
% Layout

\setlength{\parskip}{0.35em}
\setlist{itemsep=0.15em,topsep=0.3em}
\renewcommand{\topfraction}{0.9}
\renewcommand{\bottomfraction}{0.8}
\renewcommand{\textfraction}{0.07}
\renewcommand{\floatpagefraction}{0.75}
\setcounter{topnumber}{3}
\setcounter{bottomnumber}{2}
\setcounter{totalnumber}{4}
\sisetup{group-separator={\,},group-minimum-digits=5,per-mode=symbol,
  range-phrase={--},range-units=single}
\DeclareSIUnit{\cell}{cell}

% ---------------------------------------------------------------------------------------------
% Where things are, from report/ (where LaTeX runs)

\newcommand{\data}{data}        % report/data: the experiments' committed results
\newcommand{\gen}{build/gen}    % what `report.py render` writes
\newcommand{\repo}{..}          % the repository's root, to list its files
\newcommand{\repourl}{https://github.com/vehagn/vibe-rations}

% ---------------------------------------------------------------------------------------------
% Notation

\newcommand{\vect}[1]{\bm{#1}}                 % a vector: \vect{v}
\newcommand{\tens}[1]{\bm{#1}}                 % a tensor: \tens{\sigma}
\newcommand{\dt}{\Delta t}                     % the time step
\newcommand{\dx}{h}                            % the grid spacing
\newcommand{\vp}{v_p}
\newcommand{\vs}{v_s}
\newcommand{\T}{^{\mathsf T}}                  % a transpose
\newcommand{\dd}{\mathrm{d}}                   % the differential in integrals
\newcommand{\pd}[2]{\frac{\partial #1}{\partial #2}}
\DeclareMathOperator{\sinc}{sinc}

% ---------------------------------------------------------------------------------------------
% Code, files and the program's name

\newcommand{\vr}{\texttt{vibe-rations}}
% Inline code, as typed (underscores, spaces and all), breaking where a URL would:
% \code{Solver::step_with}, \code{just report-run --stale}. It prints its argument as it
% stands, macros included; for a generated value in the same type, \texttt{\val{...}}.
\DeclareUrlCommand{\vrcode}{\urlstyle{tt}}
\DeclareRobustCommand{\code}{\vrcode}
\pdfstringdefDisableCommands{\def\code#1{#1}\def\vr{vibe-rations}}
% A file of the repository, linked to it at the commit the report was built from.
\DeclareRobustCommand{\file}[1]{\href{\repourl/blob/\buildcommit/#1}{\nolinkurl{#1}}}
% A commit, shortened and linked: \commit{<the full hash, or a macro holding it>}.
\DeclareRobustCommand{\commit}[1]{\edef\vrhash{#1}\href{\repourl/commit/\vrhash}{\texttt{%
  \expandafter\shortcommit\vrhash\empty\empty\empty\empty\empty\empty\empty\empty\relax}}}
\def\shortcommit#1#2#3#4#5#6#7#8\relax{#1#2#3#4#5#6#7}

\lstdefinelanguage{toml}{
  comment=[l]{\#},
  morestring=[b]",
  morestring=[b]',
  sensitive=true,
}
\lstdefinestyle{vr}{
  basicstyle=\ttfamily\small,
  columns=fullflexible,
  keepspaces=true,
  showstringspaces=false,
  breaklines=true,
  breakatwhitespace=false,
  postbreak=\mbox{\textcolor{vrgrey}{$\hookrightarrow$}\space},
  frame=single,
  framerule=0.4pt,
  rulecolor=\color{black!25},
  backgroundcolor=\color{black!2},
  xleftmargin=4pt,
  xrightmargin=4pt,
  aboveskip=0.7em,
  belowskip=0.5em,
  commentstyle=\color{vrgrey},
  stringstyle=\color{vrgreen!70!black},
  keywordstyle=\color{vrlink},
  extendedchars=true,
  % What the solver prints and its sources' comments use beyond ASCII.
  literate=
    {×}{{$\times$}}1 {–}{{--}}1 {—}{{---}}1 {·}{{$\cdot$}}1 {−}{{$-$}}1
    {¹}{{$^1$}}1 {²}{{$^2$}}1 {³}{{$^3$}}1 {⁴}{{$^4$}}1 {⁵}{{$^5$}}1 {⁶}{{$^6$}}1
    {⁸}{{$^8$}}1 {⁻}{{$^-$}}1 {⁺}{{$^+$}}1 {⁽}{{$^($}}1 {⁾}{{$^)$}}1 {ⁿ}{{$^n$}}1
    {ᵀ}{{$^{\mathsf T}$}}1 {ᵉ}{{$^e$}}1 {ᵈ}{{$^d$}}1
    {₀}{{$_0$}}1 {₁}{{$_1$}}1 {₂}{{$_2$}}1 {ₙ}{{$_n$}}1 {ₐ}{{$_a$}}1 {ₜ}{{$_t$}}1
    {α}{{$\alpha$}}1 {γ}{{$\gamma$}}1 {δ}{{$\delta$}}1 {ε}{{$\varepsilon$}}1
    {λ}{{$\lambda$}}1 {μ}{{$\mu$}}1 {π}{{$\pi$}}1 {ρ}{{$\rho$}}1 {σ}{{$\sigma$}}1
    {τ}{{$\tau$}}1 {χ}{{$\chi$}}1 {Δ}{{$\Delta$}}1 {Σ}{{$\Sigma$}}1
    {∂}{{$\partial$}}1 {√}{{$\surd$}}1 {½}{{$\frac12$}}1 {≥}{{$\geq$}}1 {≤}{{$\leq$}}1
    {≪}{{$\ll$}}1 {≈}{{$\approx$}}1 {±}{{$\pm$}}1 {→}{{$\to$}}1 {‖}{{$\Vert$}}1
    {⌈}{{$\lceil$}}1 {⌉}{{$\rceil$}}1 {⊗}{{$\otimes$}}1 {⟨}{{$\langle$}}1
    {⟩}{{$\rangle$}}1 {ṽ}{{$\tilde v$}}1 {°}{{$^\circ$}}1 {…}{{\ldots}}1
    {’}{{'}}1 {‘}{{`}}1 {“}{{``}}1 {”}{{''}}1 {é}{{\'e}}1 {ö}{{\"o}}1 {ü}{{\"u}}1,
}
% Commands as typed; what a program printed; a file of the repository (its lines run to 100
% columns, which the smaller type fits); Rust.
\lstdefinestyle{shell}{style=vr,language=bash,keywordstyle={}}
\lstdefinestyle{output}{style=vr,basicstyle=\ttfamily\footnotesize,backgroundcolor={}}
\lstdefinestyle{toml}{style=vr,language=toml,basicstyle=\ttfamily\footnotesize}
\lstdefinestyle{rust}{style=vr,language=C,basicstyle=\ttfamily\footnotesize,
  morekeywords={fn,let,mut,pub,impl,struct,enum,match,use,mod,crate,self,Self,trait,where,
  loop,in,as,ref,move,unsafe,usize,f32,f64,bool},morecomment=[l]{///}}
\lstset{style=vr}

% A point worth remembering, set apart.
\newtcolorbox{keybox}[1][]{enhanced,breakable,colback=vrblue!4,colframe=vrlink,
  boxrule=0.5pt,arc=1.5pt,left=6pt,right=6pt,top=5pt,bottom=5pt,fonttitle=\bfseries,#1}
% Something for the reader to try.
\newtcolorbox{trybox}[1][]{enhanced,breakable,colback=vraqua!5,colframe=vraqua!60!black,
  boxrule=0.5pt,arc=1.5pt,left=6pt,right=6pt,top=5pt,bottom=5pt,fonttitle=\bfseries,
  title={Try it},#1}
% Text still to be written: visible in the PDF, so nobody takes a stub for a section.
\newcommand{\pending}[1]{\par\noindent\textcolor{vrorange!85!black}{\itshape
  [Still to be written: #1]}\par}

% ---------------------------------------------------------------------------------------------
% Figures: pgfplots reading the experiments' CSV files. The solver's curve is `solver`, an
% exact or published answer `reference`; further curves take the cycle list's colours in
% order, each with a dash pattern of its own so they survive a greyscale print.

\pgfplotsset{compat=1.18}
\pgfplotscreateplotcyclelist{vr}{
  {vrblue, solid},
  {vrorange, densely dashed},
  {vraqua!85!black, dashdotted},
  {vryellow!90!black, densely dotted},
  {vrmagenta, dashdotdotted},
  {vrgreen, loosely dashed},
  {vrviolet, solid},
  {vrred, densely dashed},
}
\pgfplotsset{
  every axis/.append style={
    width=0.82\linewidth,
    height=0.42\linewidth,
    cycle list name=vr,
    every axis plot/.append style={line width=0.8pt, mark size=1.6pt},
    grid=major,
    grid style={black!10},
    axis line style={black!55},
    tick style={black!55},
    tick label style={font=\small, /pgf/number format/fixed,
      /pgf/number format/precision=4, /pgf/number format/1000 sep={\,},
      /pgf/number format/min exponent for 1000 sep=4},
    label style={font=\small},
    title style={font=\small\bfseries},
    legend style={font=\small, draw=none, fill=white, fill opacity=0.85, text opacity=1},
    legend cell align=left,
    table/col sep=comma,
  },
  solver/.style={vrblue, solid, line width=0.9pt},
  reference/.style={black, densely dashed, line width=0.6pt},
  difference/.style={vrorange, solid, line width=0.7pt},
  % A section with z pointing down, for pictures and models: \begin{axis}[xz section, ...]
  xz section/.style={y dir=reverse, axis equal image, enlargelimits=false, grid=none,
    axis on top, xlabel={$x$ (m)}, ylabel={$z$ (m)}},
  colormap={vrdiverging}{rgb255=(27,100,190) rgb255=(239,239,239) rgb255=(214,58,59)},
  colormap={vrsequential}{rgb255=(244,248,254) rgb255=(157,200,247) rgb255=(37,125,224)
    rgb255=(14,61,126)},
}
\pgfplotstableset{
  col sep=comma,
  every head row/.style={before row=\toprule, after row=\midrule},
  every last row/.style={after row=\bottomrule},
}

% ---------------------------------------------------------------------------------------------
% What the generator provides (report/report.py render writes \gen/*.tex). Without it the
% report still compiles, with every generated value marked as missing.

\newcommand{\buildcommit}{main}
\newcommand{\buildcommitshort}{unknown}
\newcommand{\builddate}{unknown}
\newif\ifbuilddirty

\newcommand{\missing}[1]{\textcolor{vrred}{\textbf{[?\,\detokenize{#1}]}}}
\def\vrcurrent{current}
\DeclareRobustCommand{\statusmark}[1]{%
  \edef\vrstatus{#1}\ifx\vrstatus\vrcurrent current\else\textcolor{vrred}{\textbf{#1}}\fi}

% \val{usage/explosion/steps}: a value from an experiment's values.json, as it stands.
% \valnum[siunitx options]{...}: the same through \num, e.g. \valnum[round-precision=3]{...}.
% \valqty[siunitx options]{...}{unit}: with a unit.
\newcommand{\setval}[2]{\expandafter\gdef\csname vr@val@\detokenize{#1}\endcsname{#2}}
\newcommand{\val}[1]{\ifcsname vr@val@\detokenize{#1}\endcsname
  \csname vr@val@\detokenize{#1}\endcsname\else\missing{#1}\fi}
\newcommand{\valnum}[2][]{\ifcsname vr@val@\detokenize{#2}\endcsname
  \num[#1]{\csname vr@val@\detokenize{#2}\endcsname}\else\missing{#2}\fi}
\newcommand{\valqty}[3][]{\ifcsname vr@val@\detokenize{#2}\endcsname
  \qty[#1]{\csname vr@val@\detokenize{#2}\endcsname}{#3}\else\missing{#2}\fi}

% \experiment{id}{title}{tier}{status}{commit}{date}: one per manifest, from the generator.
\newcommand{\experiment}[6]{%
  \expandafter\gdef\csname vr@exp@\detokenize{#1}@title\endcsname{#2}%
  \expandafter\gdef\csname vr@exp@\detokenize{#1}@status\endcsname{#4}%
  \expandafter\gdef\csname vr@exp@\detokenize{#1}@commit\endcsname{#5}%
  \expandafter\gdef\csname vr@exp@\detokenize{#1}@date\endcsname{#6}}
% \source{usage/explosion}: where a figure's or a table's numbers come from, for the end of
% its caption: the experiment (linked to its commands in the appendix) and the commit.
\DeclareRobustCommand{\source}[1]{%
  \ifcsname vr@exp@\detokenize{#1}@status\endcsname
    \edef\vrstatus{\csname vr@exp@\detokenize{#1}@status\endcsname}%
    \edef\vrcommit{\csname vr@exp@\detokenize{#1}@commit\endcsname}%
    Data: \hyperref[exp:#1]{\texttt{\detokenize{#1}}}%
    \ifx\vrcommit\empty\else, commit \texttt{\vrcommit}\fi
    \ifx\vrstatus\vrcurrent\else\ (\statusmark{\vrstatus})\fi.%
  \else\missing{#1}\fi}
% \commands{usage/explosion}: the experiment's commands, as a listing.
\newcommand{\commands}[1]{\vrcommands#1\relax}
\def\vrcommands#1/#2\relax{\IfFileExists{\gen/commands/#1--#2.sh}
  {\lstinputlisting[style=shell]{\gen/commands/#1--#2.sh}}{\missing{#1/#2}}}
% An experiment's entry in the appendix, which \source links to.
\newcommand{\experimententry}[2]{\phantomsection\label{exp:#1}%
  \paragraph{\texttt{\detokenize{#1}}}#2.\par}

\InputIfFileExists{\gen/build.tex}{}{}
\InputIfFileExists{\gen/experiments.tex}{}{}
\InputIfFileExists{\gen/values.tex}{}{}
